\documentclass[11pt,a4paper]{article}   % Clase de documento: artículo, tamaño 11pt, papel A4
\date{}  % quita la fecha


% ===== Paquetes básicos =====
\usepackage[spanish,es-noquoting]{babel} % Soporte para español, evita comillas automáticas
\usepackage[utf8]{inputenc}              % Entrada en UTF-8 (acentos directos)
\usepackage[T1]{fontenc}                 % Codificación para salida en PDF (acentos correctos)
\usepackage{lmodern}                     % Tipografía Latin Modern (mejor que la Computer Modern)
\usepackage{geometry}                    % Márgenes personalizables
\geometry{margin=2.5cm}                  % Márgenes de 2.5 cm
\usepackage{amsmath,amssymb,amsthm}      % Paquetes AMS para matemáticas
\usepackage{hyperref}                    % Hipervínculos clickeables en índice y referencias
\hypersetup{colorlinks=true, linkcolor=blue, urlcolor=blue} % Colores de links
\usepackage{enumitem}                    % Control de espaciado en listas
\setlist{itemsep=0.2em, topsep=0.2em}    % Listas más compactas

% ===== Comandos útiles =====
\newcommand{\N}{\mathbb{N}}              % Conjunto de números naturales
\newcommand{\Pof}[1]{\mathcal{P}(#1)}    % Conjunto de partes
\newcommand{\compl}[1]{\overline{#1}}    % Complemento

% ===== Tablas y colores =====
\usepackage{array}                       % Mejor control de tablas
\newcolumntype{C}{>{$}c<{$}}             % Columna centrada en modo matemático
\usepackage{tcolorbox}                   % Cajas de texto bonitas
\usepackage[table]{xcolor}               % Colores en tablas
\usepackage{paracol}

% ===== Datos del documento =====
\title{Matemática Discreta \\[0.5em]
\large Álgebras de Boole y Relaciones de Recurrencia}

\begin{document}
\maketitle
\tableofcontents
\bigskip
\newpage

\section{Unidad 3: Álgebras de Boole}

\subsection{Conjunto de partes}
Dado un conjunto $S$, el \textit{conjunto de partes} $\Pof{S}$ es el conjunto de todos los subconjuntos de $S$.

Sea $S=\{a,b\}$. Entonces:
\[
\Pof{S}=\{\varnothing,S\}
\quad\text{(conjunto trivial)}
\]
Más explícitamente:
\[
\Pof{S}=\{\varnothing,\{a\},\{b\},S\}.
\]

\textbf{Nota.} Todos los elementos de $\Pof{S}$ se escriben entre llaves porque son conjuntos, excepto el conjunto vacío $\varnothing$.

\subsection{Álgebra de Boole de conjuntos}

\subsubsection{Relación de divisibilidad}
Si $x,y\in\N$, se dice que “$x$ divide a $y$”:
\[
x \mid y \;\Longleftrightarrow\; \exists\,k\in\N:\; y=kx.
\]

\subsubsection{Número natural libre de cuadrados}
Un número $n\in\N$ es \textit{libre de cuadrados} si ningún cuadrado mayor que $1$ lo divide.

\subsection{Álgebra de Boole de los divisores de n}
Sea un número $n\in\N$ libre de cuadrados, entonces $(D(n), \lor, \land, \lnot, 1, n)$ es un \textit{Álgebra de Boole},
donde la relación enttre sus elementos es la \textit{\textbf{divisibilidad}}.

\begin{itemize}
    \item $D(n)$ es el conjunto de todos los divisores de $n$.
    \item $x \lor y = \text{mcm}(x,y)$.
    \item $x \land y = \text{mcd}(x,y)$.
    \item $\overline x = \frac{n}{x}$.
    \item $1$ es el elemento neutro para $\lor$.
    \item $n$ es el elemento neutro para $\land$.
\end{itemize}

\subsection{Propiedades}

\begin{paracol}{2}
\paragraph{Doble complemento.}
\begin{align*}
\compl{\compl{A}}=A, \qquad
\overline{\overline{x}} = x
\end{align*}

\switchcolumn % <-- cambia a columna 2

\paragraph{Leyes de De Morgan.}
\begin{align*}
\compl{A\cup B} &= \compl{A}\cap\compl{B}, & \overline{x\lor y} &= \overline{x} \land \overline{y} \\
\compl{A\cap B} &= \compl{A}\cup\compl{B}, & \overline{x\land y} &= \overline{x} \lor \overline{y}
\end{align*}

\switchcolumn % <-- cambia a columna 1

\paragraph{Leyes conmutativas}
\begin{align*}
A \cup B &= B \cup A, & x \lor y &= y \lor x \\
A \cap B &= B \cap A, & x \land y &= y \land x
\end{align*}

\switchcolumn % <-- cambia a columna 2

\paragraph{Leyes asociativas}
\begin{align*}
A \cup (B \cup C) &= (A \cup B) \cup C, & x \lor (y \lor z) &= (x \lor y) \lor z \\
A \cap (B \cap C) &= (A \cap B) \cap C, & x \land (y \land z) &= (x \land y) \land z
\end{align*}

\switchcolumn % <-- cambia a columna 1

\paragraph{Leyes de idempotencia}
\begin{align*}
A \cup A &= A, & x \lor x &= x \\
A \cap A &= A, & x \land x &= x
\end{align*}

\switchcolumn % <-- cambia a columna 1

\paragraph{Leyes de identidad}
\begin{align*}
A \cup \varnothing &= A, & x \lor 1 &= x \\
A \cap S &= A, & x \land n &= x
\end{align*}

\switchcolumn % <-- cambia a columna 2

\paragraph{Leyes de los inversos}
\begin{align*}
A \cup \compl{A} &= S, & x \lor \overline{x} &= n \\
A \cap \compl{A} &= \varnothing, & x \land \overline{x} &= 1 \\
A \cup \compl{A} &= S, & x \lor \overline{x} &= n \\
A \cap \compl{A} &= \varnothing, & x \land \overline{x} &= 1
\end{align*}

\switchcolumn % <-- cambia a columna 1

\paragraph{Leyes de dominación}
\begin{align*}
A \cup S &= S, & x \lor n &= n \\
A \cap \varnothing &= \varnothing, & x \land 1 &= 1
\end{align*}

\switchcolumn % <-- cambia a columna 2

\paragraph{Leyes de absorción}
\begin{align*}
A \cup (A \cap B) &= A, & x \lor (x \land y) &= x \\
A \cap (A \cup B) &= A, & x \land (x \lor y) &= x
\end{align*}

\paragraph{Leyes distributivas}
\begin{align*}
A \cup (B \cap C) &= (A \cup B) \cap (A \cup C), & x \lor (y \land z) &= (x \lor y) \land (x \lor z) \\
A \cap (B \cup C) &= (A \cap B) \cup (A \cap C), & x \land (y \lor z) &= (x \land y) \lor (x \land z)
\end{align*}
\end{paracol}


\subsection{Álgebra de los interruptores}
\paragraph{Circuito en serie.} $x \cdot y = xy$
\paragraph{Circuito en paralelo.} $x + y$

\subsection{Funciones de conmutación}
Sea $B = \{0, 1\}$, el conjunto de los valores booleanos.

\paragraph{Suma} $$x + y = \begin{cases}
0 & \Leftrightarrow x = 0 \boldsymbol{\land} y = 0 \\
1 & \Leftrightarrow x = 1 \lor y = 1
\end{cases}$$

\paragraph{Producto} $$x \cdot y = \begin{cases}
0 & \Leftrightarrow x = 0 \lor y = 0 \\
1 & \Leftrightarrow x = 1 \boldsymbol{\land} y = 1
\end{cases}$$

\paragraph{Negación} $$\overline{x} = \begin{cases}
0 & \Leftrightarrow x = 1 \\
1 & \Leftrightarrow x = 0
\end{cases}$$

\subsection{Funciones booleanas}

Sea $n \in \mathbb{Z}^+$, $B^n = \{(b_1, b_2, ..., b_n) / b_i \in \{0, 1\}, \forall i\}$
($B_n$ se forma de n-uplas de 0 o 1).
Una función $f: B^n \rightarrow B$ se llama función booleana o \textit{\textbf{función de conmutación}}
y su tabla contiene $2^n$ filas.

$f:\; B^{3}\to B$, definida por $f(x,y,z)=xy+z$

\begin{center}
\begin{tabular}{C C C | C | C}
\hline
x & y & z &f(x,y,z) = xy+z \\
\hline
0 & 0 & 0 & 0\cdot 0 + 0 = 0 \\
0 & 0 & 1 & 0\cdot 0 + 1 = 1 \\
0 & 1 & 0 & 0\cdot 1 + 0 = 0 \\
0 & 1 & 1 & 0\cdot 1 + 1 = 1 \\
1 & 0 & 0 & 1\cdot 0 + 0 = 0 \\
1 & 0 & 1 & 1\cdot 0 + 1 = 1 \\
1 & 1 & 0 & 1\cdot 1 + 0 = 1 \\
1 & 1 & 1 & 1\cdot 1 + 1 = 1 \\
\hline
\end{tabular}
\end{center}

\subsection{Igualdad entre funciones booleanas}

Para $n \in \mathbb{Z}^+, n\geqslant 2$, sean $f,g:\; B^n \to B$ dos funciones booleanas, entonces
$f=g \Leftrightarrow$ las columnas de $f$ y $g$ para las mismas entradas son iguales.

\subsection{Funciones booleanas complementarias}
Para $n \in \mathbb{Z}^+, n\geqslant 2$, sea $f:\; B^n \to B$ una función booleana, entonces
$\overline{f}(x_1, ..., x_n) = \overline{f(x_1, ..., x_n)}$ 


\subsection{Adición y multiplicación de funciones booleanas}
Para $n \in \mathbb{Z}^+, n\geqslant 2$, sean $f,g:\; B^n \to B$ dos funciones booleanas, entonces:

\begin{itemize}
\item $(f+g)(x_1, ..., x_n) = f(x_1, ..., x_n) + g(x_1, ..., x_n)$
\item $(f\cdot g)(x_1, ..., x_n) = f(x_1, ..., x_n) \cdot g(x_1, ..., x_n)$
\end{itemize}

\subsection{Forma Normal Disyuntiva (FND)}
Para $n \in \mathbb{Z}^+, n\geqslant 2$, sea $f:\; B^n \to B$ una función booleana. Entonces:

\begin{enumerate}
    \item Cada $x_i$ o $\overline{x_i}$ es un \textit{\textbf{literal}}.
    \item $y_1y_2 \cdots y_n$ donde $y_i = x_i$ o $y_i = \overline{x_i}$ es una \textit{\textbf{conjunción fundamental}} (hay $2^n$ conjunciones fundamentales).
    \item $f$ está en \textit{\textbf{forma normal disyuntiva}} si es suma de conjunciones fundamentales.
\end{enumerate}

Ejemplo: $f(x,y,z) = xyz + \overline{x}y\overline{z}$ está en su \textit{\textbf{FND}}.

\begin{tcolorbox}[colback=yellow!5!white,colframe=yellow!75!black,title=Tipos de recurrencias]
Toda función booleana puede expresarse como una \textit{\textbf{adición de conjunciones fundamentales}}
y tiene su representación como una \textit{\textbf{adición de mintérminos}}.
\end{tcolorbox}


\subsection{Forma Normal Conjuntiva (FNC)}
Para $n \in \mathbb{Z}^+, n\geqslant 2$, sea $f:\; B^n \to B$ una función booleana. Entonces:

\begin{enumerate}
    \item Cada $x_i$ o $\overline{x_i}$ es un \textit{\textbf{literal}}.
    \item $y_1 + y_2 + \cdots + y_n$ donde $y_i = x_i$ o $y_i = \overline{x_i}$ es una \textit{\textbf{disyunción fundamental}}.
    \item $f$ está en \textit{\textbf{forma normal conjuntiva}} si es producto de disyunciones fundamentales.
\end{enumerate}

Ejemplo: $f(x,y,z) = (x + y + z) (x + \overline{y} + \overline{z})$ está en su \textit{\textbf{FNC}}.

\subsection{Teorema de dualidad: \texorpdfstring{$\mathrm{FND}\Leftrightarrow\mathrm{FNC}$}{FND ↔ FNC}}
\begin{itemize}
    \item Conjunciones fundamentales $\leftrightarrow$ Disyunciones fundamentales
    \item La CF se busca para $f = 1 \leftrightarrow$ la DF se busca para $f = 0$
    \item Forma normal disyuntiva $\leftrightarrow$ Forma normal conjuntiva
    \item Adición de mintérminos $\leftrightarrow$ Multiplicación de maxtérminos
\end{itemize}

\newpage
\section{Unidad 4: Relaciones de Recurencia}

\subsection{Sucesión numérica}
Una \textit{\textbf{sucesión numérica}} es una función $f:\; \N \to \mathbb{R} $.
La variable se simboliza con $n$ y la imagen de $n$ se denota $f(n) = a_n$ y se llama \textit{\textbf{término general}} de la sucesión.

\begin{tcolorbox}[colback=pink!5!white,colframe=pink!75!black,title=Recursividad]
Para construir un conjunto de forma recursiva, se necesitan:
\begin{itemize}
    \item Un conjunto inicial (o base).
    \item Una regla de formación (o paso recursivo).
\end{itemize}
\end{tcolorbox}

\subsection{Relación de recurrencia}
Una \textit{\textbf{relación de recurrencia}} es una ecuación que expresa el término general $a_n$ de una sucesión en función de uno o más términos anteriores ($a_n, a_{n-1}$, etc.).\\
Los valores necesarios para iniciar la relación se llaman \textit{\textbf{condiciones iniciales}}.

\subsection{Relaciones de recurrencia lineales}
Una relación de recurrencia es \textit{\textbf{lineal}} si el término $a_n$ es una combinación lineal de términos anteriores más una función $g(n)$ que es independiente de los términos anteriores:
\[
a_n = c_1 a_{n-1} + c_2 a_{n-2} + \cdots + c_k a_{n-k} + g(n); \forall n \geqslant k
\]

\begin{tcolorbox}[colback=yellow!5!white,colframe=yellow!75!black,title=Tipos de recurrencias]
Si $g(n) \equiv 0 $ es \textit{\textbf{homogénea}}, si no es \textit{\textbf{no homogénea}}.
\end{tcolorbox}

\subsubsection{RR lineal homogénea de orden 1}
\[
\begin{cases}
a_n = d \cdot a_{n-1}, \forall n \geqslant 1 \\
a_0 = A
\end{cases}
\]

\begin{tcolorbox}[colback=purple!5!white,colframe=purple!75!black,title=Solución]
Se define una función discreta cuyo dominio es el conjunto $\N$ de los enteros no negativos.
\[
a_n = A \cdot d^n, \forall n \geqslant 0
\]
\end{tcolorbox}

\subsubsection{RR lineal homogénea de orden 2}
Solución de forma $a_n = A \cdot r^n$ con $A \neq 0$ y $r \neq 0$
\[
\begin{cases}
C_n \cdot a_n + C_{1{n-1}} \cdot a_{n-1} + C_{2{n-1}} \cdot a_{n-2} = 0, \forall n \geqslant 2 \\
a_0; a_1
\end{cases}
\]

Y su ecuación característica es:
\[
C_n \cdot r^2 + C_1 \cdot r + C_2 = 0
\]

\begin{tcolorbox}[colback=purple!5!white,colframe=purple!75!black,title=Caso 1: Raíces reales y distintas]
\[
a_n = A_1 \cdot {r_1}^n + A_2 \cdot {r_2}^n, \forall n \geqslant 0
\]
\end{tcolorbox}

\begin{tcolorbox}[colback=purple!5!white,colframe=purple!75!black,title=Caso 2: Raíces reales e iguales]
\[
a_n = A_1 \cdot {r}^n + A_2 \cdot {r}^n \cdot n, \forall n \geqslant 0
\]
Donde se aplica el factor de corrección (el segundo término se multiplica por $n$) para que sea linealmente independiente.
\end{tcolorbox}

\begin{tcolorbox}[colback=purple!5!white,colframe=purple!75!black,title=Caso 3: Raíces complejas conjugadas]
\[
a_n = A_1 \cdot |{z}|^n \cdot \cos (n\theta)  + A_2 \cdot |{z}|^n \cdot \sin (n\theta), \forall n \geqslant 0
\]
Con $|{z}| = |{r_1}| = |{r_2}|$ y $\theta = \arg(z) = \arg(r_1) = -\arg(r_2)$, en su forma polar.
\end{tcolorbox}


\subsubsection{RR lineal no homogéneas}
\begin{tcolorbox}[colback=purple!5!white,colframe=purple!75!black,title=Método de los coeficientes indeterminados]
\[
a_n = {a_n}^{(h)} + {a_n}^{(p)}
\]
Donde:
\begin{itemize}
    \item ${a_n}^{(h)}$ es la solución de la RR \textbf{homogénea asociada}.
    \item ${a_n}^{(p)}$ es una solución particular de la RR \textbf{no homogénea}.
\end{itemize}
\end{tcolorbox}

\paragraph{Elegir ${a_n}^{(p)}$ según $g(n)$:}
\begin{center}
\begin{tabular}{|c|c|}
\hline
\rowcolor{teal!30} $g(n)$ & $a_n^{(p)}$ \\
\hline
\multicolumn{2}{|c|}{\cellcolor{teal!15} Polinómica de grado $t$} \\
$c$, cte & $A$, una constante \\
$n$ & $A_1 n + A_0$, polinomio completo de grado 1 \\
$n^2$ & $A_2 n^2 + A_1 n + A_0$, polinomio completo de grado 2 \\
$n^3$ & $A_3 n^3 + A_2 n^2 + A_1 n + A_0$, polinomio completo de grado 3 \\
$n^t$ & polinomio completo de grado $t$ \\
\hline
\multicolumn{2}{|c|}{\cellcolor{teal!15} Exponencial de base $r$} \\
$r^n$ & $A r^n$, una constante por la misma exponencial \\
\hline
\multicolumn{2}{|c|}{\cellcolor{teal!15} Trigonométricas} \\
$\sen(\theta n)$ & $A_1 \sen(\theta n) + A_2 \cos(\theta n)$ \\
$\cos(\theta n)$ & $A_1 \sen(\theta n) + A_2 \cos(\theta n)$ \\
\hline
\end{tabular}
\end{center}

\begin{tcolorbox}[colback=yellow!5!white,colframe=yellow!75!black,title=Recordar]
Si la forma propuesta para ${a_n}^{(p)}$ entra en conflicto con algún término de la solución de la RR homogénea asociada, se multiplica por $n^k$,
donde $k$ es el menor entero positivo tal que la nueva forma propuesta no sea solución de la RR homogénea asociada.
\end{tcolorbox}

\fbox {\textbf{Nota:} en algún punto de la resolución (casi) siempre los términos más grandes de $n$ se anulan.}

\end{document}
